% !TEX root = ../main.tex
% !TEX encoding = UTF-8 Unicode
% !TEX TS-program = pdflatexmk
% !TEX spellcheck = en-US
% !BIB program = biber

%%%%% CHAPTER 1
% want quotes?
%\begin{savequote}[8cm]
%Computers are incredibly fast, accurate and stupid; humans are incredibly slow, inaccurate, and brilliant; together they are powerful beyond imagination
%
%\qauthor{--- Albert Einstein}
%
%\end{savequote}

\chapter{Introduction}\label{chap:intro}%
%

% want an abstract per chapter??
\begin{chapterabstract}
	This chapter introduces the main ideas that went into this thesis. The first section gives a historically inspired motivation, beginning with the inception of quantum field theory, specifically quantum electrodynamics and quantum chromodynamics. We then explain how the understanding of strong-coupling phenomena in quantum field theories motivates the introduction of supersymmetry, which allowed Seiberg and Witten to draw a connection to integrable models. This thesis is introduced as a natural extension of these developments. The second section recounts the temporal progression of ideas that led to the results of this thesis. Since the author comes from a mathematical background, our exposition will be targeted towards mathematicians.
\end{chapterabstract}

% want a toc per chapter??
%\minitoc

\section{Motivation}\label{sec:motivation}

Two great revolutions in physics---quantum mechanics and relativity---occurred at the beginning of the 20th century. On one hand, quantum mechanics describes observables in the world by hermitian linear operators and symmetries of the world by unitary linear operators, both acting on a Hilbert space of physical states. On the other hand, relativity describes observables in the world by sections of bundles over a four-dimensional spacetime manifold $M$ equipped with a metric with Lorentzian signature $(+,-,-,-)$, and the symmetries of the world are described by the Lorentz group $\mathrm{SO}(1,3)$ (or the simply connected double cover $\mathrm{Spin}(1,3)$ of its identity component) of linear transformations preserving the metric and its extension by spacetime translations $P_\mu$ for coordinate indices $\mu=0,1,2,3$, the Poincaré group $\mathrm{SO}(1,3) \ltimes \R^4$.

Restricting to Minkowski spacetime $M=\R^{1,3}$ with the standard Minkowski metric $\eta = \operatorname{diag}(+1,-1,-1,-1)$, meaning
\begin{equation}
	(ds)^2 = \eta_{\mu\nu} dx^\mu dx^\nu = (dx^0)^2 - (dx^1)^2 - (dx^2)^2 - (dx^3)^2,
\end{equation}
where repeated indices are understood to be summed over, the approaches of quantum mechanics and relativity are unified in the framework of quantum field theory \cite{peskin_introduction_2018}, which posits that the Hilbert space of physical states and the algebra of observables acting on it come equipped with compatible unitary actions of the Poincaré group. Observables $\Phi$ acting on the Hilbert space then become \emph{fields}, \emph{i.e.}\ dependent on a position $x = (x^0,x^1,x^2,x^3)$ in spacetime, via
\begin{equation}
	\Phi(x) \coloneq e^{\I x^\mu P_\mu} \Phi e^{-\I x^\mu P_\mu}.
\end{equation}
Rigorously speaking, fields should be interpreted as operator-valued distributions, which may be integrated over regions in spacetime. It then becomes relevant to ask how the fields transform under the action of the Lorentz group. Due to the isomorphism
\begin{equation}
	\mathfrak{spin}(1,3)_\C \cong \mathfrak{sl}(2,\C) \oplus \mathfrak{sl}(2,\C)
\end{equation}
of complex Lie algebras, representations can be described by two half-integers $(j_L,j_R)$ which are the spin labels or weights of finite-dimensional irreducible representations of the left- and right-handed $\mathfrak{sl}(2,\C)$. Physically relevant fields usually transform in the scalar representation $(0,0)$, the left- or right-handed spinor representations $(\tfrac 12,0)$ or $(0,\tfrac 12)$, or the vector representation $(\tfrac 12, \tfrac 12)$.

A simple example of a quantum field theory is the free massive scalar quantum field theory, which describes a real valued field $\Phi(x)$ (think: hermitian operator-valued distribution) that transforms in the scalar representation $(0,0)$ of the Lorentz group and whose dynamics is described by the action functional
\begin{equation}
	S_\text{free}[\Phi] \coloneq \int_{\R^{1,3}} d^4 x \left[ \frac{1}{2} \partial^\mu \Phi \partial_\mu \Phi - \frac{m^2}{2} \Phi^2 \right]
\end{equation}
for some mass $m > 0$, where the spacetime index is raised via the inverse metric $\eta^{\mu\nu} = \eta_{\mu\nu}$, \emph{i.e.} $\partial^\mu = \eta^{\mu\nu} \partial_\nu$. This means that, classically, the field $\Phi$ must extremize the action functional,
\begin{equation}
	\frac{\delta S_\text{free}}{\delta \Phi} = 0,
\end{equation}
and that, quantum mechanically, observables $\mathcal{O}$ (diagonal in the field basis) have expectation values
\begin{equation}
	\langle \mathcal{O} \rangle \coloneq Z^{-1} \int [d\Phi] e^{\tfrac{\I}{\hbar} S_\text{free}[\Phi]} \mathcal{O}[\Phi], \qquad Z \coloneq \int [d\Phi] e^{\tfrac{\I}{\hbar} S_\text{free}[\Phi]},
\end{equation}
where $Z$ is the partition function and $[d\Phi]$ is the notoriously elusive path integral measure. In the particular case of the free massive scalar, however, expectation values can be calculated explicitly, making the free massive scalar quantum field theory an \emph{exactly solvable} quantum field theory.

A more involved example of a quantum field theory that is also relevant to the real world is \emph{quantum electrodynamics}, describing the interactions of electrons, positrons, and photons, which is also our first example of a \emph{gauge theory}. The theory includes two fields: The \emph{gauge field} $A(x) = A_\mu(x) dx^\mu$ transforming in the vector representation $(\tfrac 12, \tfrac 12)$---the photon field; and the \emph{Dirac spinor field} $\psi(x)$ transforming in the Dirac spinor representation $(\tfrac 12,0) \oplus (0,\tfrac 12)$---the electron-positron field. What makes this quantum field theory a gauge theory is that the fields transform under an additional symmetry, called \emph{gauge symmetry}. Specifically, smooth maps $h\colon \R^{1,3} \to U(1)$ generate the transformations
\begin{equation}
	A_\mu(x) \mapsto h(x) A_\mu(x) h(x)^{-1} + e^{-1} h(x) \partial_\mu h(x)^{-1}, \qquad \psi(x) \mapsto h(x) \psi(x),
\end{equation}
where $e$ is the coupling constant controlling the strength of electromagnetic interactions. This exhibits the gauge field $A(x)$ as a principal $U(1)$-connection 1-form for the trivial principal $U(1)$-bundle over $\R^{1,3}$. The dynamics of quantum electrodynamics are then described by the essentially unique gauge invariant action functional
\begin{equation}
	S_\text{QED}[A,\psi] \coloneq \int_{\R^{1,3}} d^4 x \left[ {-\frac{1}{4}} F^{\mu\nu} F_{\mu\nu} + \bar \psi (\I \gamma^\mu D_\mu - m) \psi \right],
\end{equation}
with the electron mass $m > 0$, the curvature 2-form $F \coloneq dA$, the covariant derivative
\begin{equation}
	D_\mu \coloneq \partial_\mu + e A_\mu,
\end{equation}
and the Dirac matrices $\gamma^\mu$, which map the tensor product of the Dirac spinor representation with its dual to the vector representation.

Despite the mathematically beautiful differential-geometric formulation of quantum electrodynamics, the interactions induced by the presence of the term $e \bar \psi \gamma^\mu A_\mu \psi$ in the action place closed form expectation values out of reach of analytical methods. Even perturbative methods, where $e$ is assumed to be small, are famously plagued by divergent integrals. Perturbative expectation values can nonetheless be extracted by \emph{renormalization}, which regularizes the divergent integrals by cutting off the integration range at an energy scale $k$ and introducing $k$-dependent counter terms that modify the values of the coupling constant $e$ and the electron mass $m$ in a way that cancels the divergent terms. For this reason, one says that both the coupling and the mass must \emph{run} with the energy scale $k$. Importantly, the coupling constant $e$ must become small at low energy scales, which makes the perturbative predictions of quantum electrodynamics (coming in the form of asymptotic series in $e$) valid at ordinary energy scales.

However, not every physically relevant gauge theory has a valid perturbative description at low energy scales. In fact, \emph{quantum chromodynamics}, describing the interactions of quarks and gluons relevant in nuclear physics, has the opposite behavior: Its coupling constant becomes small at high energies---a property known as \emph{asymptotic freedom}---and large at low energies \cite{gross_ultraviolet_1973}. This invalidates perturbative methods for calculations at ordinary energy scales, requiring theoreticians to develop non-perturbative methods valid in the strong coupling regime. Everyday experience and numerical lattice simulations suggest that strongly coupled quantum chromodynamics exhibits a phenomenon known as \emph{confinement}, binding charged quarks into neutral composite configurations known as hadrons.

To further explore strongly coupled phenomena in quantum field theories, a simplified testing ground is instructive. The most basic example of an asymptotically free quantum field theory is Yang--Mills theory \cite{yang_conservation_1954} with gauge group $G_c$, describing the dynamics of a gauge field $A(x)$ valued in $\mathfrak{g} \coloneq \operatorname{Lie} G_c$ (think: principal connection 1-form) transforming in the vector representation $(\tfrac 12, \tfrac 12)$ of the Lorentz group with the action functional
\begin{equation}
	S_\text{YM}[A] \coloneq \int_{\R^{1,3}} d^4 x \left[ {-\frac{1}{4}} \Tr F^{\mu\nu} F_{\mu\nu} \right], \qquad F \coloneq dA + g A \wedge A,
\end{equation}
where $g$ is the coupling constant, determining the strength of non-abelian self-interaction of the gauge field. Indeed, Yang--Mills theory for $G_c = \mathrm{SU}(N)$ faces the same problems as quantum chromodynamics: The weak coupling approximation becomes invalid at low energies.

Surprisingly, confinement can be seen semiclassically after placing the theory on the spacetime $\R^{1,2} \times S_R^1$ with small circle radius $R$ and adding additional fields to Yang--Mills theory \cite{poppitz_continuity_2012}. Specifically, adding massless $\mathfrak{g}$-valued fields $\lambda(x),\bar \lambda(x)$ transforming in the \emph{Majorana} spinor representation of the Poincaré algebra leads to the action functional
\begin{equation}
	S_{\text{SYM}}[A,\lambda] \coloneq \int_{\R^{1,3}} d^4 x \Tr \left[ {-\frac{1}{4}} F^{\mu\nu} F_{\mu\nu} - \lambda \I \sigma^\mu D_\mu \bar \lambda \right],
\end{equation}
where $\sigma^\mu$ are the Pauli matrices, which map the tensor product of the Majorana spinor representation with its dual to the vector representation. This quantum field theory is known as $\mathcal{N}=1$ supersymmetric Yang--Mills theory, since it possesses an extended symmetry which extends the Poincaré algebra to a Lie superalgebra using additional generators of odd degree, called \emph{supercharges}, $Q$ and $\bar Q$, which transform in the $(\tfrac 12,0)$ and $(0,\tfrac 12)$ representations under the Lorentz group, and importantly satisfy
\begin{equation}
	[Q,\bar Q]_+ \coloneq Q \bar Q + \bar Q Q = 2\sigma^\mu P_\mu.
\end{equation}
This extended Lie superalgebra is known as the $\mathcal{N}=1$ super Poincaré algebra, capturing the spacetime symmetries of super Minkowski spacetime $\R^{1,3|4}$. Using the formalism of supermanifolds, the fields $A(x)$ and $\lambda(x),\bar \lambda(x)$ combine to transform in the \emph{vector multiplet representation} of the $\mathcal{N}=1$ super Poincaré algebra.

Continuing in this direction, a revolutionary breakthrough was achieved by Seiberg, Witten, and others after further extending supersymmetry to $\mathcal{N}=2$ supersymmetric Yang--Mills theory with gauge group $\mathrm{SU}(N)$ \cite{seiberg_monopole_1994,seiberg_monopoles_1994,gorsky_integrability_1995,martinec_integrable_1996,donagi_supersymmetric_1996,dhoker_calogero-moser_1998}. This quantum field theory extends $\mathcal{N}=1$ supersymmetric Yang--Mills theory by another set of massless $\mathfrak{su}(N)$-valued Majorana spinor fields $\psi(x),\bar \psi(x)$ as well as a massless $\mathfrak{su}(N)_\C$-valued scalar field $\varphi(x) = \phi_1(x) + \I \phi_2(x)$. All of these fields may be combined into a field $\Psi(x,\theta,\bar \theta)$ on super Minkowski spacetime $\R^{1,3|8}$, which transforms in the \emph{vector multiplet representation} of the $\mathcal{N}=2$ super Poincaré algebra, capturing the spacetime symmetries of $\R^{1,3|8}$. The action functional may be conveniently packaged as an integral over one chiral half of super Minkowski spacetime
\begin{equation}\label{eq:prepotentialAction}
	S_{\text{SYM2}}[\Psi] \coloneq \operatorname{Im} \frac{1}{4\pi} \int d^4 x \, d^4 \theta \left[ \frac{1}{4} \Tr \mathcal{F}(\Psi) \right],
\end{equation}
where the \emph{prepotential} $\mathcal{F}$ is some holomorphic function. At high energies, the theory has the prepotential $\mathcal{F}(\Psi) \coloneq \frac{\tau}{2} \Psi^2$ with $\tau = \frac{\theta}{2\pi} + \frac{4\pi\I}{g^2}$, where $g$ is the coupling constant and $\theta$ is an angular variable called the \emph{$\theta$-angle}.

The important new ingredient present in $\mathcal{N}=2$ super Yang--Mills theory is the adjoint-valued complex scalar field $\varphi(x) = \phi_1(x) + \I \phi_2(x)$. At minimum energy, the potential forces
\begin{equation}
	[\phi_1(x),\phi_2(x)] = 0,
\end{equation}
so that the expectation values $\langle \varphi \rangle$ can be taken to lie in the diagonal Cartan subalgebra of $\mathfrak{su}(N)_\C$, whose components $a^i$ (modulo the Weyl group $S_N$) coordinatize the moduli space of Coulomb vacua, the \emph{Coulomb branch} $\mathfrak{B}$, as a complex variety. Expanding the field $\varphi$ in the action functional \eqref{eq:prepotentialAction} around the point $\langle \varphi \rangle$, one can see that all off-diagonal components of the $\mathfrak{su}(N)$-valued fields acquire a non-zero mass proportional to $a^i-a^j$. At low energy scales, we may integrate over all massive modes so that only the massless fields survive in the effective theory. It follows that all fields are diagonal in the resulting \emph{effective theory}, which is characterized by the values of the gauge invariant expectation values $u = (u_2,\dots,u_N)$ with $u_k = \langle \Tr \varphi^k \rangle = \sum_{i=1}^N (a^i)^k$. This simplifies the effective description enormously.

Crucially, supersymmetry tells us that the actional functional of the effective theory must still be of the form \eqref{eq:prepotentialAction}, albeit with a possibly different prepotential $\mathcal{F}$. Expanding the prepotential to quadratic order, the effective coupling matrix has components
\begin{equation}
	\tau_{ij} = \frac{\partial^2 \mathcal{F}}{\partial a^i \partial a^j}
\end{equation}
and $\operatorname{Im} \tau$ must be positive definite to ensure that the effective theory is unitary. This equips the Coulomb branch $\mathfrak{B}$ with a special Kähler structure where the Kähler metric is given by
\begin{equation}
	(ds)^2 = (\operatorname{Im} \tau)_{ij} da^i d\bar a^j.
\end{equation}
The effective theory can thus be written as a four-dimensional non-linear sigma model on $\R^{1,3}$ whose target space is the Coulomb branch $\mathfrak{B}$ with this Kähler metric.

Arguably, the main motivation of Seiberg and Witten to do this construction was that it allows to show the existence of a confining phase in a mass deformation of the theory. Indeed, one can use the algebraic relations of the $\mathcal{N}=2$ super Poincaré algebra to show that the mass of particles of electric charge $n$ and magnetic charge $m$ must satisfy the BPS bound \cite{witten_supersymmetry_1978}
\begin{equation}
	M \geq \sqrt{2} |n \cdot a + m \cdot a_D|, \qquad a_D \coloneq \frac{\partial \mathcal{F}}{\partial a}.
\end{equation}
Particles that saturate this bound are called \emph{BPS particles}. Usually, BPS particles of non-zero charge are heavy and thus do not contribute to the dynamics of the effective theory at low energy scales. However, it can happen at particular points of the Coulomb branch $\mathfrak{B}$ that $a \to 0$ or $a_D \to 0$. In this case, certain BPS particles become light and cannot be integrated out of the effective theory. Turning on a non-zero mass for the coordinate $u_2$ (breaking supersymmetry to $\mathcal{N}=1$) to drive the theory towards the points $a \to 0$ or $a_D \to 0$ also drives the light BPS particle to a non-zero expectation value---it \emph{condenses}. If this BPS particle is a magnetic monopole (carrying no electric charge), the monopole condensate increases the energy of electric flux times area. The electric flux therefore gets squeezed into tubes with a small cross section which have a \emph{string tension}, energetically favoring shorter flux tubes. It follows that sources of (color-)electric charge connected by electric flux tubes are strongly attracted and bind into (color-)neutral configurations, \emph{i.e.\ }become \emph{confined}.

In order for the description of the light BPS monopoles to be valid perturbatively, we need the magnetic interactions to be weak. It thus becomes necessary to study the behavior of the coupling matrix $\tau$ more closely. Importantly, $\tau$ defines a local system of lattices $\Gamma$, where the lattice over a point $u$ is given by
\begin{equation}
	\Gamma_u \coloneq \Z^{N-1} + \tau(u) \Z^{N-1}.
\end{equation}
This, in turn, defines a torus fibration $\mathfrak{M} \coloneq T^*\mathfrak{B}/\Gamma \to \mathfrak{B}$ over the Coulomb branch $\mathfrak{B}$, which comes with the canonical holomorphic symplectic form of the cotangent bundle. \emph{The punchline is that the coordinates $u$ form a Poisson-commutative family turning the torus fibration $\mathfrak{M}$ into the phase space of an algebraic integrable system in the sense of the Liouville theorem.} Specifically, it turns out that $a^i$ and their canonical conjugates are the action and angle variables of the periodic Toda chain (in the center-of-mass frame), whose Hamiltonian may be written as
\begin{equation}
	u_2 = H_\text{Toda} = \frac{1}{2} \sum_{i=1}^{N} p_i^2 + \Lambda^2 \sum_{i=1}^{N} e^{x_i - x_{i+1}},
\end{equation}
where $x_i$ are particle coordinates, $p_i$ the canonical momenta, and $i$ is taken modulo $N$. This Hamiltonian has the stationary points
\begin{equation}
	p_i = 0, \qquad x_i - x_{i+1} = 2\pi\I\,k/N,
\end{equation}
at which $u_2 = N\Lambda^2 e^{2\pi\I\,k/N}$. Since these points are stationary, the Liouville torus at these energy levels (the fiber sitting over a point $u$) must degenerate in the direction of the flow. The modulus $\tau$ of the torus therefore blows up in this direction, making the effective theory of the corresponding light BPS particle weakly coupled.

A very general family of 4d $\mathcal{N}=2$ gauge theories (referred to as class $\mathcal{S}$ theories) can be constructed by compactifying the 6d $\mathcal{N}=(2,0)$ superconformal field theory of type $A_{N-1}$ on a punctured Riemann surface $C$. In this case, \cite{gaiotto_wall-crossing_2011} showed that the analogous torus fibration $\mathfrak{M} \to \mathfrak{B}$ should be identified with the \emph{Hitchin fibration}, whose total space $\mathfrak{M}$ is the moduli space of Higgs bundles and whose base $\mathfrak{B}$ (the Coulomb branch of the theory on $\R^{1,3}$) is given by the Hitchin base
\begin{equation}
	\mathfrak{B} = \bigoplus_{k=2}^N H^0(C,K_C^{\otimes k}),
\end{equation}
where $K_C$ is the canonical bundle of $C$. The corresponding integrable model is the \emph{Hitchin system} associated to the punctured Riemann surface $C$.

To further explore the integrable structure of the Coulomb branch, we can again compactify the spacetime to $\R^{1,2} \times S_R^1$ \cite{seiberg_gauge_1996,gaiotto_four-dimensional_2010,gaiotto_wall-crossing_2011,gaiotto_three_2013}. According to the Kaluza--Klein construction, only the Fourier zero modes of the fields along the circle are massless. All other modes are massive and do not contribute to the dynamics of the effective theory on $\R^{1,2}$. However, the compactification introduces new scalars to the effective theory: The zero mode $\phi_3$ of the gauge field component $A_3$ along the circle and the dual photon $\sigma$ parametrizing the three-dimensional Hodge dual of the zero mode $F_3$ of the curvature via
\begin{equation}
	\star_3 F_3 = d\sigma.
\end{equation}
Large gauge transformations send $\phi_3$ and $\sigma$ to
\begin{equation}\label{eq:scalarTransLaw}
	\phi_3 \mapsto \phi_3 + \Z^{N-1}, \qquad \sigma \mapsto \sigma + \tau \Z^{N-1}
\end{equation}
Therefore, the gauge invariant combination of $\phi_3$ and $\sigma$ is the monopole operator
\begin{equation}\label{eq:monoOp}
	v \coloneq e^{\frac{2\pi\I}{\tau} (\tau \phi_3 + \sigma)}.
\end{equation}
The transformation law \eqref{eq:scalarTransLaw} says that the monopole operator $v$ is a complex coordinate on the torus fibers $T_u^* \mathfrak{B}/\Gamma_u$. The special Kähler structure on $\mathfrak{B}$ may then be extended to a hyperkähler structure on $\mathfrak{M}$ with the \emph{semiflat metric}
\begin{equation}
	(ds)^2 = R (\operatorname{Im} \tau)_{ij} da^i d\bar a^j + R^{-1} (\operatorname{Im} \tau)^{-1}_{ij} \frac{dv^i}{v^i} \frac{d\bar v^j}{\bar v^j}.
\end{equation}

The sigma model on $\R^{1,2}$ with target space given by $\mathfrak{M}$ equipped with the semiflat metric approximates the theory on $\R^{1,2} \times S_R^1$, but crucially misses the dynamics of BPS particles with worldlines wrapping around the circle. Taking these dynamics into account, the complex symplectic variety obtained from the small circle limit $R \to 0$ (which is independent of the choice of complex structure) was closely considered in \cite{bullimore_coulomb_2015} and rigorously defined by Braverman, Finkelberg, Nakajima, and others (BFN) in a seminal series of papers \cite{nakajima_towards_2016,braverman_towards_2018,braverman_coulomb_2018}. This construction generalizes to the case of general gauge groups $G_c$ and additional ``matter'' fields transforming in a cotangent-type representation $T^* V$ under the gauge group and in the hypermultiplet representation under the $\mathcal{N}=2$ super Poincaré algebra.

Looking instead at the holomorphic symplectic form of the complex structure $J$ at finite radius $R$ (as opposed to the complex structure $I$ coming from the coordinates $u$ and their canonical conjugates), the Darboux charts form cluster charts in the sense of cluster varieties, where the Poisson bracket of the Darboux coordinates $X_\gamma$ is
\begin{equation}
	\{ X_\gamma, X_{\gamma'} \} = \langle \gamma,\gamma' \rangle X_{\gamma+\gamma'}.
\end{equation}
Following Braverman--Finkelberg--Nakajima, the resulting complex symplectic variety also has a mathematically rigorous construction for general gauge groups $G_c$ and hypermultiplet fields in cotangent-type representations $T^* V$, known as \emph{$K$-theoretic Coulomb branches}. Their cluster structure was conjectured in \cite{gaiotto_four-dimensional_2010,gaiotto_three_2013} and found by Schrader and Shapiro \cite{schrader_k-theoretic_2026}. Generally, $K$-theoretic Coulomb branches also furnish algebraic integrable systems. Note, however, that the complex structure $J$ is distinct from the complex structure $I$, and the case $G_c=\mathrm{SU}(N)$ and $V=0$ does not yield the periodic Toda chain as in the complex structure $I$, but the open relativistic Toda chain. For class $\mathcal{S}$ theories, the complex structure $J$ yields the character variety of the Riemann surface $C$, instead of the moduli space of Higgs bundles in complex structure $I$.

Following \cite{nekrasov_five_1998}, the idea of Seiberg and Witten can be repeated for five-dimension\-al gauge theories with $\mathcal{N}=1$ supersymmetry placed on $\R^{1,2} \times T^2$: The effective theory can be rewritten as a sigma model on $\R^{1,2}$ with values in the Coulomb branch $\mathfrak{M}$ of the theory, which is again a hyperkähler manifold whose metric now depends on the radii $R_3$ and $R_4$ of $T^2$. The holomorphic symplectic forms of the complex structures $I$ and $J$ again define complex symplectic varieties. For example, in the case $G_c=\mathrm{SU}(N),V=0$, the complex structure $I$ gives rise to the periodic \emph{relativistic} Toda chain. The complex symplectic variety coming from the complex structure $J$ again fits (conjecturally) into the BFN framework for any $G_c$ and $V$ using equivariant elliptic cohomology. 

To summarize, the correspondence between supersymmetric gauge theories with gauge group $\mathrm{SU}(N)$ and integrable models living on the Coulomb branch is given in table \ref{table:intSysPure} without matter, and in table \ref{table:intSysAdjoint} with adjoint matter, which is the class $\mathcal{S}$ theory for the once-punctured torus. The bottom right entries are extrapolations based on the rational--trigonometric--elliptic hierarchy of the momentum dependence and the entries of the first row may be obtained from the entries of the second row by degeneration limits.
\begin{table}[h]
	\centering
	\begin{tabular}{|l|c|c|}
		\hline
		$G_c=\mathrm{SU}(N),V=0$ & complex structure $I$ & complex structure $J$ \\
		\hline
		$\mathcal{N}=4, \R^{1,2}$ & open Toda & open Toda \\
		$\mathcal{N}=2, \R^{1,2} \times S^1$ & periodic Toda \cite{gorsky_integrability_1995,martinec_integrable_1996} & open relativistic Toda \cite{bezrukavnikov_equivariant_2026} \\
		$\mathcal{N}=1, \R^{1,2} \times T^2$ & periodic relativistic Toda \cite{nekrasov_five_1998} & dual periodic Toda? \\ \hline
	\end{tabular}
	\caption[Coulomb branches without matter]{Integrable models living on the Coulomb branches of gauge theories with gauge group $G_c = \mathrm{SU}(N)$ without matter.\label{table:intSysPure}}
\end{table}
\begin{table}[h]
	\centering
	\begin{tabular}{|l|c|c|}
		\hline
		$G_c=\mathrm{SU}(N),V=\mathfrak{su}(N)$ & complex structure $I$ & complex structure $J$ \\
		\hline
		$\mathcal{N}=4, \R^{1,2}$ & trigonometric CM & trigonometric CM \\
		$\mathcal{N}=2, \R^{1,2} \times S^1$ & elliptic CM \cite{dhoker_calogero-moser_1998} & trigonometric RS \cite{braverman_cyclotomic_2020} \\
		$\mathcal{N}=1, \R^{1,2} \times T^2$ & elliptic RS \cite{nekrasov_five_1998} & dual elliptic RS? \\ \hline
\end{tabular}
	\caption[Coulomb branches with adjoint matter]{Integrable models living on the Coulomb branches of gauge theories with gauge group $G_c = \mathrm{SU}(N)$ with adjoint matter.\label{table:intSysAdjoint}}
\end{table}

The main results of this thesis can be understood as a natural continuation of the ideas above to the case of the \emph{necklace quiver gauge theory}. The gauge group and matter content of the necklace quiver gauge theory can be read off of the necklace quiver shown in figure \ref{fig:quiver}. For every node $\alpha \in \Z_\ell$ in the quiver (drawn as circles filled with the number $N$), we get a factor $\mathrm{U}(N)_\alpha$ in the gauge group (each associated to a $\mathfrak{u}(N)$-valued vector multiplet), and for every arrow, we get a hypermultiplet valued in the cotangent representation of the bifundamental representation $V_\alpha \otimes V_{\alpha+1}^\vee$, where $V_\alpha$ is the defining representation of $\mathrm{U}(N)_\alpha$ and $V_\alpha^\vee$ is its dual. This means
\begin{equation}
    G_c = \prod_{\alpha \in \Z_\ell} \mathrm{U}(N)_\alpha, \qquad V = \bigoplus_{\alpha \in \Z_\ell} V_\alpha \otimes V_{\alpha+1}^\vee.
\end{equation}
\begin{figure}[h]
    \centering
    \def\l{6}
    \def\radius{1.2cm}
    \begin{tikzpicture}[>=Stealth, every node/.style={font=\small}]
        \node[circle,draw,inner sep=1.5pt] (v1) at ({360*(0)/\l}:\radius) {$N$};
        \node[circle,draw,inner sep=1.5pt] (v2) at ({360*(1)/\l}:\radius) {$N$};
        \node[circle,draw,inner sep=1.5pt] (v3) at ({360*(2)/\l}:\radius) {$N$};
        \node[circle,draw,inner sep=1.5pt] (v4) at ({360*(3)/\l}:\radius) {$N$};
        \node (v5) at ({360*(4)/\l}:\radius) {$\cdots$};
        \node[circle,draw,inner sep=1.5pt] (v6) at ({360*(5)/\l}:\radius) {$N$};
        
        \foreach \i in {1,...,\l} {
            \pgfmathtruncatemacro{\next}{int(mod(\i,\l)+1)}
            \draw[->, shorten >=2pt, shorten <=2pt] (v\i) to[bend right=14] (v\next);
        }
        
        \node at ({(180+360*(0))/\l}:{1.2*\radius}) {\small $m^1$};
        \node at ({(180+360*(1))/\l}:{1.2*\radius}) {\small $m^2$};
        \node at ({(180+360*(2))/\l}:{1.2*\radius}) {\small $m^3$};
        \node at ({(180+360*(3))/\l}:{1.2*\radius}) {\small $m^4$};
        \node at ({(180+360*(4))/\l}:{1.2*\radius}) {\small $m^{\ell-1}$};
        \node at ({(180+360*(5))/\l}:{1.2*\radius}) {\small $m^0$};
    \end{tikzpicture}
    \caption[The necklace quiver]{The necklace quiver with $\ell$ nodes of rank $N$. The bifundamental hypermultiplets corresponding to the arrows carry masses $m^\alpha$.}
    \label{fig:quiver}
\end{figure}
The results of this thesis show that the Coulomb branches of necklace quiver gauge theories with $\ell$ nodes of rank $N$ are the phase spaces of the spin Ruijsenaars--Schneider models of $N$ particles with $\ell$ spin states. However, care should be taken, as the Hamiltonians defining the spin Ruijsenaars--Schneider models are actually \emph{spectrally dual} to the Hamiltonians Seiberg--Witten theory would dictate. The precise meaning of this statement will be explained in the main text.
\begin{table}[h]
	\centering
	\begin{tabular}{|l|c|c|}
		\hline
		Necklace quiver & complex structure $J$ \\
		\hline
		$\mathcal{N}=4, \R^{1,2}$ & dual rational spin RS \\
		$\mathcal{N}=2, \R^{1,2} \times S^1$ & dual trigonometric spin RS \\
		$\mathcal{N}=1, \R^{1,2} \times T^2$ & dual elliptic spin RS \\ \hline
\end{tabular}
	\caption[Coulomb branches of necklace quiver gauge theories]{Integrable models living on the Coulomb branches of necklace quiver gauge theories.\label{table:intSysNecklace}}
\end{table}

Table~\ref{table:intSysNecklace} exhibits the dictionary that the remainder of this thesis makes precise. Reading down the column, the compactification data of the spacetime fixes the degeneration type of the integrable system: for $\mathbb{R}^{1,2}$ the model is rational, for $\mathbb{R}^{1,2}\times S^1$ the model is trigonometric, and for $\mathbb{R}^{1,2} \times T^2$ the model is elliptic, meaning that the particles move on $\C$, $\C^\times$, or an elliptic curve, respectively. This is the same hierarchy that relates ordinary cohomology, $K$-theory and elliptic cohomology, which is precisely the input datum distinguishing the three versions of the Braverman--Finkelberg--Nakajima construction. Read off from the quiver of figure~\ref{fig:quiver}, the gauge-theoretic data fixes the kinematics of the spin Ruijsenaars--Schneider model: the rank $N$ of the gauge group at each node is the number of particles, and the number $\ell$ of nodes is the number of spin states each particle carries. The momenta, by contrast, are multiplicative in all three rows, since in the complex structure $J$ the fiber coordinates are the monopole operators $v$ of \eqref{eq:monoOp}. This is why every entry of table~\ref{table:intSysNecklace} is of Ruijsenaars--Schneider rather than Calogero--Moser type.

What this identification buys is uniformity. Poisson structures for spin Ruijsenaars--Schneider models have so far been produced one degeneration at a time and by different means: by Hamiltonian reduction in the rational case~\cite{arutyunov_hamiltonian_1998}, by quasi-Hamiltonian reduction in the trigonometric case, and in the elliptic case only on the spectral data of the Lax matrix~\cite{krichever_elliptic_1998}, with no expression in the natural variables $x_i$, $a^\alpha_i$, $c^\alpha_i$ beyond $N=2$~\cite{soloviev_hamiltonian_2008}. The Coulomb branch construction is indifferent to this distinction: one and the same quiver, evaluated in the appropriate cohomology theory, yields the Poisson algebra together with a distinguished commuting family of Hamiltonians, and its quantization is obtained by replacing the Coulomb branch with its quantized counterpart.

\section{Progression of the doctorate}

As often happens in research, the starting point for this thesis was of a very different nature than what was laid out in the previous section. The research question this thesis set out to solve was to write down a Poisson algebra, Hamiltonian, and quantization for the equations of motion of the elliptic spin Ruijsenaars--Schneider integrable model found by Krichever and Zabrodin \cite{krichever_spin_1995}. We recount them here for the readers convenience:
\begin{equation}\label{eq:eoms}
    \begin{aligned}
        \dot x_i
        ={}& {-f_{ii}}, \\[6pt]
        \dot a_i^\alpha
        ={}& \sum_{j=1}^N (V(x_i-x_j) f_{ij} a_i^\alpha - V(x_i-x_j) f_{ij} a_j^\alpha), \\[-1pt]
        \dot c_i^\alpha
        ={}& \sum_{j=1}^N (V(x_j-x_i) f_{ji} c_j^\alpha - V(x_i-x_j) f_{ij} c_i^\alpha).
    \end{aligned}
\end{equation}
These equations of motion describe the dynamics of $N$ particles with $\ell$ spin states. The variables $x_i$ with particle index $i=1,\dots,N$ denote particle positions, and the variables $a_i^\alpha$ and $c_i^\alpha$ denote the components of the spin covector and spin vector, respectively, where $\alpha \in \Z_\ell$ is the spin state index. The inner product between the spin covectors and spin vectors is abbreviated by
\begin{equation}
	f_{ij} \coloneq \sum_{\alpha \in \Z_\ell} a_i^\alpha c_j^\alpha
\end{equation}
and the potential is
\begin{equation}
	V(z) \coloneq E_1(z) - E_1(z-\eta),
\end{equation}
where $E_1(z)$ is the logarithmic derivative of the odd Jacobi theta function $\theta(z)$ depending on the modulus $\tau$ of the elliptic curve
\begin{equation}
	E_\tau \coloneq \frac{\C}{\Z+\tau \Z}
\end{equation}
on which the particles move. The equations of motion possess a symmetry which allow a simultaneous rescaling of the spin covector and spin vector:
\begin{equation}
	a_i^\alpha \mapsto k_i a_i^\alpha, \qquad c_i^\alpha \mapsto k_i^{-1} c_i^\alpha.
\end{equation}
This symmetry can be fixed by imposing the constraint
\begin{equation}
	a_i^0 = 1.
\end{equation}
Furthermore, the equations of motion are symmetric under the action of the spin rotation group $\mathrm{GL}_\ell$ by acting on the spin covectors and spin vectors as covectors and vectors, respectively. The constraint $a_i^0 = 1$ is preserved by the mirabolic subgroup $\mathrm{M}_\ell$ of $\mathrm{GL}_\ell$. The invariant dynamical variables of the elliptic spin Ruijsenaars--Schneider model therefore respect the constraint $a_i^0 = 1$ and must be $\mathrm{M}_\ell$-invariant combinations of the variables $x_i,a_i^\alpha,c_i^\alpha$. Before taking $\mathrm{M}_\ell$-invariants, the phase space has dimension $2N\ell$, while after taking invariants, the resulting phase space has dimension $2N\ell - \ell(\ell-1)$. Note that this is in fact an even number, allowing for the existence of a symplectic form. In particular, in the case $\ell = 1$, the only dynamical variables are $x_i$ and $c_i^0$ for $i=1,\dots,N$, whose equations of motion reduce to the equations of motion of the elliptic (spinless) Ruijsenaars--Schneider model \cite{ruijsenaars_complete_1987}, also known as the relativistic elliptic Calogero--Moser model.

Krichever and Zabrodin showed that the spin Ruijsenaars--Schneider model is an integrable model by exhibiting a Lax pair. This means that there are two matrices $L(z)$ and $M(z)$ satisfying the Lax equation
\begin{equation}
	\dot L(z) = [M(z),L(z)].
\end{equation}
Specifically, the Lax pair of Krichever and Zabrodin is given by
\begin{equation}
    \begin{aligned}
		L(z) \coloneq{}& \sum_{i,j=1}^N \Phi(z,x_j-x_i+\eta) f_{ij} e_{ij}, \qquad \\
        M(z) \coloneq{}& \sum_{i=1}^N E_1(z) f_{ii} e_{ii} + \sum_{i,j=1}^N V(x_i-x_j) f_{ij} e_{ii} + \sum_{i \neq j}^N \Phi(z,x_j-x_i) f_{ij} e_{ij}.
    \end{aligned}
\end{equation}
Subsequently, Krichever constructed a symplectic form for the $\mathrm{M}_\ell$-invariant variables from the spectral data of the Lax matrix \cite{krichever_elliptic_1998}. However, it was not immediately obvious how to reexpress this symplectic form in terms of the natural variables $x_i,a_i^\alpha,c_i^\alpha$, except in the spinless case $\ell=1$. In the case $N=2$, \cite{soloviev_hamiltonian_2008} wrote down a symplectic form for the variables
\begin{equation}
	x_1, x_2, f_{11}, f_{22},
\end{equation}
where the dynamics of the variables $f_{12},f_{21}$ are integrated out and $\eta$ is replaced by an effective coupling $z_0$ defined by the equation $\det L(z_0) = 0$.

Constructing Poisson structures in more general cases required degenerating the underlying elliptic curve, yielding the trigonometric and rational spin Ruijsenaars--Schneider models. In the rational case, where the elliptic curve degenerates to a plane, a Poisson structure was found by Hamiltonian reduction \cite{arutyunov_hamiltonian_1998}. There, it was shown that the rational spin Ruijsenaars--Schneider model of $N$ particles and $\ell$ spin states can be obtained as a Hamiltonian reduction
\begin{equation}\label{eq:ratRSPhaseSpace}
	(T^* \mathrm{GL}_N \times (T^* \C^N)^{\times \ell}) \hred \mathrm{GL}_N,
\end{equation}
where $\mathrm{GL}_N$ acts on $\mathrm{GL}_N$ by the adjoint action and on $\C^N$ by the defining representation. Via the multiplication map on $\mathrm{GL}_N$, we can trivialize the cotangent bundle, giving an isomorphism
\begin{equation}
	T^* \mathrm{GL}_N \cong \mathrm{GL}_N \times \mathfrak{gl}_N^* \cong \mathrm{GL}_N \times \mathfrak{gl}_N,
\end{equation}
where we used the trace inner product for the last isomorphism. We thus obtain a parametrization of $T^* \mathrm{GL}_N$ by $g \in \mathrm{GL}_N$ and $x \in \mathfrak{gl}_N$, and of $(T^* \C^N)^{\times \ell}$ by $a \in \C^{N \times \ell}$ and $c \in \C^{\ell \times N}$. The action of $\mathrm{GL}_N$ can then be expressed as
\begin{equation}\label{eq:adjAction}
	g \mapsto hgh^{-1}, \qquad x \mapsto hxh^{-1}, \qquad a \mapsto ha, \qquad c \mapsto c h^{-1}.
\end{equation}
Again identifying $\mathfrak{gl}_N^* \cong \mathfrak{gl}_N$, the momentum map $\mu\colon T^* \mathrm{GL}_N \times (T^* \C^N)^{\times \ell} \to \mathfrak{gl}_N^*$ for the action \eqref{eq:adjAction} turns out to be
\begin{equation}
	\mu = x - gxg^{-1} - ac.
\end{equation}
The phase space \eqref{eq:ratRSPhaseSpace} of the (spinless) rational Ruijsenaars--Schneider model is obtained by imposing the momentum map equation $\mu = \eta \cdot \id$, where $\eta$ is the coupling constant, and quotienting by the action of $\mathrm{GL}_N$. The Poisson-commuting Hamiltonians of the Ruijsenaars--Schneider model are given by $H[n] = \Tr g^n$. In particular, the Hamiltonian $H[1]$ generates the equations of motion \eqref{eq:eoms}.

Given this construction of the phase space of the rational spin Ruijsenaars--Schneider model by Hamiltonian reduction, chapter \ref{chap:spinRSRed} quantizes the model using the method of quantum Hamiltonian reduction. The main motivation for carefully quantizing the rational case was the observation that it already contained an unexplored structure that develops its full traction in the quantization. Namely, besides finding a Poisson algebra that reproduces the equations of motion \eqref{eq:eoms}, the paper \cite{arutyunov_hamiltonian_1998} also identified two important symmetry algebras inside the Poisson algebra generated by the combinations
\begin{equation}
	T(z) \coloneq 1 + c (z-x)^{-1} a, \qquad J[n] \coloneq c g^n a,
\end{equation}
whose Poisson brackets are of the form
\begin{align}
	\{ T_1(z),T_2(w) \} &= \frac{1}{z-w} [P,T_1(z) T_2(w)], \\
	\{ J_1[n],J_2[m] \} &= \frac{1}{2} [P,J_1[n+m]-J_2[n+m]],
\end{align}
where $P = \sum_{\alpha,\beta \in \Z_\ell} e_{\alpha\beta} \otimes e_{\beta\alpha}$ is the permutation operator on $\C^\ell \otimes \C^\ell$. Now, the important insight is that these Poisson brackets are the defining Poisson brackets of the classical limit of the Yangian $Y(\mathfrak{gl}_\ell)$ and the loop algebra $L(\mathfrak{gl}_\ell)$. More precisely, $T(z)$ and $J[n]$ generate the $N$-truncated quotients thereof, which comes from the fact that $T(z)$ and $J[n]$ are defined in terms of $N \times N$ matrices. Moreover, we find that
\begin{equation}\label{eq:HamFromCurrent}
	\Tr J[n] = \Tr c g^n a = \Tr ac g^n = -\eta H[n],
\end{equation}
where we used the momentum map equation for the last equality. The Hamiltonians of the rational spin Ruijsenaars--Schneider model therefore come from the \emph{center} of the loop algebra $L(\mathfrak{gl}_\ell)$.

At this point, it was natural to conjecture that the $N$-truncated Yangian and loop algebra combine into a larger algebra which is the $N$-truncated loop Yangian $LY(\mathfrak{gl}_\ell)$. This conjecture was also motivated by the construction of other spin generalizations of Ruijsenaars--Schneider Hamiltonians from the quantum toroidal algebra \cite{lamers_spin-ruijsenaars_2022}, whose rational degeneration is the loop Yangian. Through many conversations, it became known to the author that truncated shifted Yangians are examples of Coulomb branches of three-dimensional quantum field theories with $\mathcal{N}=4$ supersymmetry. The natural continuation of the search for a Poisson structure of the elliptic spin Ruijsenaars--Schneider model thus was to have a closer look at Coulomb branches.

The paper \cite{bullimore_coulomb_2015} offered very concrete insight into the Poisson structures defined by Coulomb branches. In particular, section 6 included a description of the momentum map of the enhanced topological symmetry $\mathfrak{sl}_\ell$ acting on the Coulomb branch of linear quiver gauge theories. Importantly, the $(\alpha,\beta)$-coefficients of the momentum map were expressed as sums of the form
\begin{equation}
	\sum_{i^\alpha,i^{\alpha+1},\dots,i^\beta} \frac{u_{i^\alpha}^\alpha \cdots u_{i^\beta}^\beta}{(\varphi_{i^{\alpha+1}}^{\alpha+1}-\varphi_{i^\alpha}^\alpha) \cdots (\varphi_{i^\beta}^{\beta}-\varphi_{i^{\beta-1}}^{\beta-1})}
\end{equation}
that could be written as matrix products of $L$-operators of the form
\begin{equation}
	L^\alpha = \sum_{i,j=1}^N \frac{u_j^{\alpha+1}}{\varphi_j^{\alpha+1} - \varphi_i^\alpha} e_{ij},
\end{equation}
resembling the Lax matrix of the (spinless) Ruijsenaars--Schneider model. Comparing these formulas to the ideas of \cite{nakajima_cherkis_2025} made it clear that such a construction should also be possible in the case of the necklace quiver. Constructed correctly, the momentum map for the necklace quiver then takes values in the loop algebra $L(\mathfrak{gl}_\ell)$, and the trace of the first mode should give the defining Hamiltonian $H[1]$ of the rational spin Ruijsenaars--Schneider model as in equation \eqref{eq:HamFromCurrent}. From the knowledge of the Hamiltonian $H[1]$ and the identification $x_i = \varphi_i^0$, we may compute $f_{ii}$ from the equations of motion \eqref{eq:eoms}, which inspires an ansatz for the spin (co)vector variables $a_i^\alpha$ and $c_i^\alpha$. Again using the equations of motion as a check, this ansatz finally yielded the correct identification of variables between the Coulomb branch variables and the variables of the rational spin Ruijsenaars--Schneider model.

After finding the correct identification of variables in the rational case, the next step was the extension to the trigonometric and elliptic cases, corresponding to the Coulomb branches constructed from equivariant $K$-theory and equivariant elliptic Borel--Moore homology instead of equivariant Borel--Moore homology. The extension to the trigonometric case followed swiftly, but the elliptic case remained resistant due to the necessity of a spectral parameter, which in the rational and trigonometric cases may simply be taken to infinity using the non-compactness of $\C$ and $\C^\times$. Only after a lot of careful matching to the references \cite{arutyunov_r-matrix_1997,maruyoshi_wilson-t_2021} could the identification of variables be found in the elliptic case. This identification contained a big surprise: The Poisson structure underlying the elliptic model includes dynamical inhomogeneities without which the Jacobi identity fails. Arguably, the apparent necessity of dynamical inhomogeneities is the crux of why the Poisson structure of the elliptic model has eluded us for over three decades. We hope that the requirement for dynamical inhomogeneities can be lifted in the future and a rigorous comparison with Krichever's symplectic structure can be made.

\section{Structure of the thesis}